\documentclass[aspectratio=169]{beamer}
\usepackage{will_handley_beamer}
\usepackage{title_page}

% Commands
% --------
% - \arxiv{arxiv number}
% - \arxiv{<number>}            arxiv.org/abs/<number>
% - \oldarxiv{<arxiv number>}   arxiv.org/<number>
% - \doi{<doi>}                 doi.org/<doi>
% - \xkcd{<number>}             xkcd.com/<number>
% - \email{<email>}             <<email>>
% - \tthref{<website>}          <website>
% - \av[dist]{<quantity>}       <quantity>_{dist}
% - \student{<name>}{<detail>}{<photo>}

% Talk details
% ------------
\title{Cosmic Tensions: A Particle Physicist's Primer}
\date{1\textsuperscript{st} November 2024}

\begin{document}

\begin{frame}
    \titlepage
\end{frame}

\begin{frame}
    \frametitle{Introduction: Precision Cosmology and its Discontents}
    \begin{columns}
        \column{0.5\textwidth}
        \begin{itemize}
            \item Cosmology has well-and-truly entered an era of precise measurements, analogous to particle physics.
            \item Multiple independent observations allow us to constrain the parameters of our cosmological model.
            \item The Standard Model of Cosmology ($\Lambda$CDM) successfully explains a wide range of observations.
            \item However, increasing precision has revealed inconsistencies, or tensions, between different measurements.
            \item Are these tensions cracks in $\Lambda$CDM, hints of new physics, or simply measurement systematics?
        \end{itemize}
        \column{0.5\textwidth}
        \includegraphics[width=\textwidth]{figures/Logo_rgb.pdf}
    \end{columns}
\end{frame}


\begin{frame}
    \frametitle{Measuring Distances in the Universe: A Cosmological Toolkit}
    \begin{columns}
        \column{0.5\textwidth}
        \begin{itemize}
            \item \textbf{Standard Candles:} Objects with known intrinsic luminosity (e.g., Type Ia supernovae).  Measuring their apparent brightness allows us to infer their distance. \arxiv{2112.04510}
            \item \textbf{Standard Rulers:} Objects with known physical size (e.g., Baryon Acoustic Oscillations). Measuring their apparent size allows us to infer their distance. \arxiv{2404.03000}
            \item \textbf{Standard Sirens:} Merging compact objects (e.g., neutron stars, black holes). The gravitational wave signal provides a direct measurement of their distance. \arxiv{1710.05835}
        \end{itemize}
        \column{0.5\textwidth}
        \includegraphics[width=\textwidth]{figures/Dist_Ladd_distance_ladder_R22.pdf}
    \end{columns}
\end{frame}


\begin{frame}
    \frametitle{The Hubble Constant: A Tale of Two Values}
    \begin{columns}
        \column{0.5\textwidth}
        \begin{itemize}
            \item $H_0$ quantifies the current expansion rate of the Universe.
            \item Early Universe probes (CMB) infer $H_0 \approx 67$ \kms. \arxiv{1807.06209}
            \item Late Universe probes (distance ladder) measure $H_0 \approx 73$ \kms. \arxiv{2112.04510}
            \item This $\sim 5\sigma$ discrepancy is the Hubble tension.
        \end{itemize}
        \column{0.5\textwidth}
        \includegraphics[width=\textwidth]{figures/Intro_H0_whisk.pdf}
    \end{columns}
\end{frame}

\begin{frame}
    \frametitle{The Sound Horizon: Another Perspective on Expansion}
    \begin{columns}
        \column{0.5\textwidth}
        \begin{itemize}
            \item The sound horizon ($r_d$) is the distance sound waves travelled in the early Universe. It acts as a standard ruler.
            \item CMB infers $r_d \approx 147$ Mpc. \arxiv{1807.06209}
            \item Late-time measurements suggest $r_d \approx 137$ Mpc.
            \item This discrepancy is linked to the Hubble tension.
        \end{itemize}
        \column{0.5\textwidth}
        \includegraphics[width=\textwidth]{figures/BAO_CosmoVerse_HubbleDiag_2D_final.pdf}
    \end{columns}
\end{frame}

\begin{frame}
    \frametitle{The Clustering of Matter: A Smoother Universe?}
    \begin{columns}
        \column{0.5\textwidth}
        \begin{itemize}
            \item $S_8$ quantifies the amplitude of matter fluctuations on large scales.
            \item CMB infers a larger $S_8$ than weak lensing surveys. This is the $S_8$ tension.
            \item Could indicate less structure in the late Universe than predicted.
        \end{itemize}
        \column{0.5\textwidth}
        \includegraphics[width=\textwidth]{figures/Intro_S8_whisk.pdf}
    \end{columns}
\end{frame}

\begin{frame}
    \frametitle{Cosmology vs. Particle Physics: A Tale of Two Disciplines}
    \begin{columns}
        \column{0.5\textwidth}
        \begin{itemize}
            \item Cosmology: Single realization of the Universe.  Observations are limited by cosmic variance.
            \item Particle Physics: Multiple, repeatable experiments. Statistical uncertainties dominate.
            \item These differences necessitate different statistical approaches.
        \end{itemize}
        \column{0.5\textwidth}
        \includegraphics[width=\textwidth]{figures/cross_correlation_schematic_1.pdf}
    \end{columns}
\end{frame}

\begin{frame}
    \frametitle{Bayesian vs. Frequentist Inference}
    \begin{columns}
        \column{0.5\textwidth}
        \begin{itemize}
            \item \textbf{Bayesian:}  Incorporates prior knowledge.  Provides probability distributions for parameters.
            \item \textbf{Frequentist:} Relies solely on data.  Provides confidence intervals for parameters.
            \item Both approaches have their pros and cons in cosmology.
        \end{itemize}
        \column{0.5\textwidth}
        \includegraphics[width=\textwidth]{figures/BAO_CosmoVerse_fig_3a_final_rectangle.pdf}
    \end{columns}
\end{frame}

\begin{frame}
    \frametitle{Marginalization vs. Profiling: A Matter of Perspective}
    \begin{columns}
        \column{0.5\textwidth}
        \begin{itemize}
            \item \textbf{Marginalization:} Integrating over nuisance parameters. Common in Bayesian analyses.  Can be affected by volume effects.
            \item \textbf{Profiling:} Minimizing over nuisance parameters. Common in frequentist analyses. Less sensitive to priors.
            \item Understanding these differences is crucial for interpreting results. \arxiv{2408.07700}
        \end{itemize}
        \column{0.5\textwidth}
        \includegraphics[width=\textwidth]{figures/PMF_MCMC.pdf}
    \end{columns}
\end{frame}


\begin{frame}
    \frametitle{Nuisance Parameters: The Unseen Hand}
    \begin{columns}
        \column{0.5\textwidth}
        \begin{itemize}
            \item Parameters that affect the data but are not of primary interest (e.g., dust extinction, calibration uncertainties).
            \item Can introduce degeneracies and bias cosmological parameter estimates.
            \item Careful treatment is essential for robust results.
        \end{itemize}
        \column{0.5\textwidth}
        \includegraphics[width=\textwidth]{figures/snia_H0_syst.pdf}
    \end{columns}
\end{frame}


\begin{frame}
    \frametitle{The Dark Energy Tension: Hints of Dynamics?}
    \begin{columns}
        \column{0.5\textwidth}
        \begin{itemize}
            \item The equation of state of dark energy, $w$, describes its properties.
            \item $\Lambda$CDM assumes $w = -1$ (cosmological constant).
            \item Some data suggests $w \neq -1$, hinting at dynamic dark energy.  \arxiv{2307.14802}
            \item The Dark Energy Spectroscopic Instrument (DESI) will provide crucial data.
        \end{itemize}
        \column{0.5\textwidth}
        \includegraphics[width=\textwidth]{figures/Whisker_plot_w.pdf}
    \end{columns}
\end{frame}

\begin{frame}
    \frametitle{The Future of Cosmic Tensions}
    \begin{columns}
        \column{0.5\textwidth}
        \begin{itemize}
            \item New surveys (Euclid, Rubin Observatory LSST, Roman Space Telescope) will provide unprecedented data.
            \item Advanced data analysis techniques (simulation-based inference, machine learning) will be crucial. \arxiv{2401.17029}
            \item Removing fiducial cosmology assumptions will improve robustness.
            \item Lensed calibrators will provide independent distance measurements.
        \end{itemize}
        \column{0.5\textwidth}
        \includegraphics[width=\textwidth]{figures/Mod_Grav.pdf}
    \end{columns}
\end{frame}

\begin{frame}
    \frametitle{Conclusion: A Call to Arms}
    \begin{columns}
        \column{0.5\textwidth}
        \begin{itemize}
            \item Cosmic tensions are a major challenge to our understanding of the Universe.
            \item Could lead to revolutionary discoveries in cosmology and fundamental physics.
            \item Interdisciplinary collaborations between particle physicists, astrophysicists, and cosmologists will be essential.
        \end{itemize}
        \column{0.5\textwidth}
        \includegraphics[width=\textwidth]{figures/IntroFunSols.pdf}
    \end{columns}
\end{frame}

\begin{frame}
    \frametitle{Conclusions}
    \framesubtitle{\tthref{github.com/handley-lab}}
    \tikz[overlay,remember picture]
        \node[anchor=north east] (A) at ($(current page.north east)+(0,0)$) {
        \includegraphics[width=0.10\textheight]{people/adam_ormondroyd.jpg}%
        \includegraphics[width=0.10\textheight]{people/charlotte_priestley.jpg}%
        \includegraphics[width=0.10\textheight]{people/david_yallup.jpg}%
        \includegraphics[width=0.10\textheight]{people/dily_ong.jpg}%
        \includegraphics[width=0.10\textheight]{people/george_carter.jpg}%
        \includegraphics[width=0.10\textheight]{people/harry_bevins.jpg}%
        \includegraphics[width=0.10\textheight]{people/kilian_scheutwinkel.jpg}%
        \includegraphics[width=0.10\textheight]{people/metha_prathaban.jpg}%
        \includegraphics[width=0.10\textheight]{people/namu_kroupa.jpg}%
        \includegraphics[width=0.10\textheight]{people/sinah_legner.jpg}%
        \includegraphics[width=0.10\textheight]{people/sam_leeney.jpg}%
        \includegraphics[width=0.10\textheight]{people/wei-ning_deng.jpg}%
        \includegraphics[width=0.10\textheight]{people/will_handley.jpg}%
    };
\end{frame}


\end{document}
